\section{Related Work}
\label{sec:related}

\subsection{Dough mechanics and inverse rheological modeling}

Food-engineering studies provide both precedents for inverse simulation and reasons to limit the interpretation of a fitted robot model. Fabbri and Cevoli identify the consistency and flow indices of a power-law semolina-dough model by inverting finite-element extrusion simulations against force-derived pressure measurements~\cite{fabbri2015}. Their comparison with capillary rheometry evaluates a characterization procedure, rather than prediction from partial visual observations. Cocuzza and Yan identify a spring--dashpot model of fondant icing using controlled mechanical tests before evaluating gravity- and tool-induced sheet deformation~\cite{cocuzza2018}. These studies show that inverse modeling of food materials predates recent differentiable robotic simulators, and that the selected constitutive equation determines what its parameters mean.

Flour-dough rheology includes nonlinear and history-dependent effects that are not fully represented by a compact elastic--viscous--plastic approximation~\cite{yazar2023}. Residual deformation after a short recording does not uniquely demonstrate yield plasticity: slow recovery, viscous flow, and interface adhesion can produce similar observations. Adhesive behavior also depends on probe material and test conditions~\cite{adhesion2014}. A principal-stretch limit should therefore be reported as a dimensionless model parameter, not renamed a measured yield stress. For the present study, controlled loading rates, unloading, and recovery observations are important complements to complex two-tool motion.

\subsection{Learned models for elastoplastic manipulation}

RoboCraft reconstructs particle representations from RGB-D observations and learns graph-network dynamics for elastoplastic object manipulation~\cite{robocraft2022}. RoboCook extends learned prediction and planning to diverse tools and long-horizon manipulation of real dough~\cite{robocook2023}. These results establish real food manipulation as an existing robotics application, rather than an unexplored use of deformable simulation. DoughNet predicts geometry and topological changes from a single RGB-D observation and tool actions~\cite{doughnet2024}, demonstrating that limited-view action-conditioned prediction is also studied outside physical parameter identification.

Learned dynamics and compact constitutive identification address related but distinct goals. A predictive network can be useful without assigning a mechanical interpretation to its latent representation. Conversely, explicit coefficients are not automatically more accurate or identifiable. A comparison should therefore match the claim: forward prediction under the same recorded actions, data required for calibration, or consistency with independent mechanical tests. Final target-reaching scores from a planner are not directly comparable to geometric errors from a prescribed-motion replay.

\subsection{Elastoplastic material point identification}

DPSI is a close precedent for the proposed inverse problem~\cite{dpsi2025}. It combines fixed-corotated elasticity, von Mises plasticity, MLS-MPM/APIC transfer, and rigid-tool signed-distance contact. A robot executes known trajectories on plasticine; one camera is moved to six viewpoints before and after manipulation, and the endpoint clouds are fused. Thus, one physical camera supplies multi-view endpoints rather than a fixed-view time series. The method fits material parameters and, in a higher-contact-complexity setting, interface friction. It also evaluates unseen, longer tool motions. Known actions, incomplete-volume completion, elastoplastic fitting, and action-transfer evaluation are therefore established contributions. Its reported parameter ambiguities and model-dependent recovery behavior further motivate evaluating prediction separately from parameter uniqueness.

EMPM extends embodied material point modeling to offline demonstrations and online bimanual interaction, explicitly including bread dough and observations from three RGB-D cameras~\cite{empm2026}. Its general formulation includes elastic and plastic quantities, while the reported experimental implementation optimizes homogeneous Young's modulus and Poisson's ratio. Gradient-based fitting and a forward-simulator CMA-ES option are both described. The online bread-dough experiment evaluates alignment with ongoing parameter correction; this differs from fixing parameters before an independent manipulation. Our proposed comparison concerns the value of time-resolved, single-view evidence and a separately defined rate term, not the novelty of two arms or dough simulation. The principal-stretch projection used here also differs from the von Mises laws in DPSI and EMPM and must be evaluated on its own merits.

\subsection{Visual identification of elastic and dissipative response}

Arnavaz et al. fit soft-body simulations directly to depth images from a single camera~\cite{diffcal2023}. With known silicone-specimen geometry and clamp conditions, their dynamic example jointly estimates Young's modulus, a damping coefficient, and density. They explicitly discuss dependence on the constitutive model and discretization. Their damping coefficient is not interchangeable with the additive viscosity coefficient used here. This work establishes single-depth-view elastic and dissipative calibration, while leaving different measurement and contact assumptions from an initially incompletely observed dough object.

DiffCloud combines differentiable simulation and rendering with two-camera point clouds and recorded robot motion for cloth and other deformable-object experiments~\cite{diffcloud2022}. Its real fitting examples use selected informative frames and a directed simulated-to-observed point-distance loss. The present directed term has the opposite orientation, from observations to visible predicted particles, and is combined with time-resolved depth and coverage penalties. The choice of direction determines which unsupported regions receive a penalty; neither directed distance is equivalent to a symmetric Chamfer loss.

Earlier inverse mechanics studies also inform the evaluation. Wang et al. estimate soft-object models from tracked multi-view motion and include independent loading comparisons~\cite{wang2015}. Hahn et al. use marker motion capture and measured clamp motion to identify elastic and power-law viscous response, showing the importance of varied excitation and boundary conditions~\cite{hahn2019}. These studies support testing multiple motions and numerical resolutions rather than interpreting optimizer convergence as sufficient evidence of material recovery.

\subsection{Visual continuum models and interaction-conditioned prediction}

PAC-NeRF couples differentiable material point simulation to radiance-field reconstruction for multi-view system identification~\cite{pacnerf2023}. GIC uses Gaussian-informed geometry and temporal visual objectives for continuum-parameter identification and simulation~\cite{gic2024}. Both consider multiple constitutive families in synthetic experiments, including dissipative or plastic behavior. Their broader formulation should be distinguished from the particular real experiments that were demonstrated; coverage of a constitutive family in synthetic data does not establish joint rheological identification of manipulated flour dough.

PhysTwin provides a complementary spring--mass approach fitted to multi-view hand--object interactions~\cite{phystwin2025}. It evaluates unseen interactions by reusing the fitted model after registration to the new episode's first observation. This is a relevant precedent for separating material reuse from initial-state reconstruction. In our proposed evaluation, a held-out episode likewise receives its own first-frame geometry and recorded actions, but its subsequent dough observations must not influence calibration or parameter selection.

\subsection{Numerical foundations and the role of gradients}

MLS-MPM supplies the particle--grid discretization~\cite{mls2018}, while APIC preserves a locally affine velocity representation during transfer~\cite{apic2015}. Stomakhin et al.'s snow model provides a precedent for principal-stretch plastic projection with corotated elasticity~\cite{snow2013}; it does not validate that constitutive family for flour dough. ChainQueen, DiffTaichi, and PlasticineLab establish differentiable material point simulation, differentiable programming, and elastoplastic manipulation benchmarks~\cite{chainqueen2019,difftaichi2020,plasticinelab2021}. Reusing these numerical ideas is a methodological choice, not by itself a research contribution.

Gradients can improve optimization efficiency but do not supply missing experimental information. Different parameter combinations can remain practically indistinguishable under partial observation~\cite{raue2009}, and a fitted parameter can compensate for systematic model discrepancy~\cite{discrepancy2014}. The scientific question is therefore whether the proposed measurements constrain useful predictions across motions, with uncertainty and numerical dependence reported. Endpoint versus time-resolved fitting, restricted versus joint parameter sets, and gradient versus derivative-free search under comparable budgets are appropriate tests of that question.
